% ======================================================================
%  PART 10 -- Non-Linear Time Series: ARCH and GARCH
% ======================================================================
\section{Non-Linear Time Series: ARCH and GARCH}
\label{sec:archgarch}

Linear models (AR, MA, ARMA) describe the conditional \emph{mean} of a series and
assume a constant innovation variance. Financial log-returns violate this: they
are nearly uncorrelated (so a linear conditional-mean model is almost useless),
yet their \emph{squares} are strongly dependent --- large moves cluster with
large moves and quiet periods with quiet periods. This \textbf{volatility
clustering}, together with heavy (\textbf{fat}) tails, motivates models for the
\emph{conditional variance}. \textbf{ARCH} and \textbf{GARCH} keep the series
mean-zero and uncorrelated (white noise) while letting the conditional variance
$\sigma_t^2$ evolve in time as a function of past squared returns (and, for
GARCH, past variances). Throughout, $\mathcal{F}_t$ denotes the information
available up to and including time $t$, and $\{\eps_t\}$ is white noise with
$\E[\eps_t]=0$, $\Var(\eps_t)=1$, independent of $\mathcal{F}_{t-1}$.

% ----------------------------------------------------------------------
\subsection{Definitions}

\begin{definition}{Log-returns \& volatility clustering}{p10-logret}
For a price series $\{X_t\}$ (e.g.\ a stock index), the \textbf{log-returns} are
\[
r_t=\log\!\Big(\frac{X_t}{X_{t-1}}\Big)=\log(X_t)-\log(X_{t-1}),
\]
measuring relative changes in price. Empirically log-returns show:
\begin{itemize}
  \item \textbf{little linear dependence:} the sample ACF/PACF of $r_t$ is
        essentially that of white noise;
  \item \textbf{dependent squares:} the sample ACF/PACF of $r_t^2$ decays slowly
        $\Rightarrow$ the conditional variance changes over time;
  \item \textbf{volatility clustering:} large changes tend to be followed by
        large changes, small by small;
  \item \textbf{fat tails:} kurtosis $>3$ (heavier tails than the normal),
        visible as deviations in a QQ-plot against $\mathcal{N}(0,1)$.
\end{itemize}
\end{definition}

\begin{definition}{ARCH(1) model}{p10-arch1}
$\{r_t\}$ is an \textbf{AutoRegressive Conditionally Heteroscedastic} model of
order $1$, $\mathrm{ARCH}(1)$, if for $t\in\Ints$
\[
r_t=\sigma_t\eps_t,\qquad \sigma_t^2=\alpha_0+\alpha_1 r_{t-1}^2,
\]
with $\{\eps_t\}$ white noise of mean $0$ and variance $1$, independent of
$\mathcal F_{t-1}$, and parameters $\alpha_0>0,\ \alpha_1\ge 0$. The volatility
$\sigma_t$ is $\mathcal F_{t-1}$-measurable (known one step ahead). Immediate
properties (Proposition~\ref{prop:p10-archmoments}):
\[
\E[r_t\mid\mathcal F_{t-1}]=0,\qquad
\Var(r_t\mid\mathcal F_{t-1})=\sigma_t^2=\alpha_0+\alpha_1 r_{t-1}^2 .
\]
\end{definition}

\begin{definition}{ARCH(p) model}{p10-archp}
$\{r_t\}$ is an $\mathrm{ARCH}(p)$ model if $r_t=\sigma_t\eps_t$ with
\[
\sigma_t^2=\alpha_0+\alpha_1 r_{t-1}^2+\cdots+\alpha_p r_{t-p}^2
=\alpha_0+\sum_{j=1}^{p}\alpha_j r_{t-j}^2,
\]
$\{\eps_t\}$ white noise $(0,1)$ independent of $\mathcal F_{t-1}$, and
$\alpha_0>0,\ \alpha_p>0,\ \alpha_j\ge0$ for $1\le j<p$. Volatility may now
depend on several past squared shocks, so as before
$\Var(r_t\mid\mathcal F_{t-1})=\sigma_t^2$.
\end{definition}

\begin{definition}{GARCH(p,q) model}{p10-garch}
$\{r_t\}$ is a \textbf{Generalised ARCH} model $\mathrm{GARCH}(p,q)$ if
$r_t=\sigma_t\eps_t$ with
\[
\sigma_t^2=\alpha_0+\sum_{j=1}^{p}\alpha_j r_{t-j}^2+\sum_{j=1}^{q}\beta_j\sigma_{t-j}^2,
\]
$\{\eps_t\}$ white noise $(0,1)$ independent of $\mathcal F_{t-1}$, and parameters
$\alpha_0>0$, $\alpha_p>0$, $\beta_q>0$, with $\alpha_j,\beta_j\ge0$ otherwise.
Feeding back past \emph{variances} $\sigma_{t-j}^2$ lets volatility persist over
time far more parsimoniously than an $\mathrm{ARCH}(p)$ with large $p$. A
$\mathrm{GARCH}(p,0)$ is exactly an $\mathrm{ARCH}(p)$. As before
$\Var(r_t\mid\mathcal F_{t-1})=\sigma_t^2$.
\end{definition}

\begin{definition}{Standardised residuals (diagnostics)}{p10-resid}
After fitting, the \textbf{standardised residuals} are
$\tilde e_t=r_t/\hat\sigma_t$. If the model is correct these should look like
i.i.d.\ $(0,1)$ noise: their ACF/PACF \emph{and} the ACF/PACF of their
\emph{squares} $\tilde e_t^2$ should resemble white noise, and a QQ-plot against
the assumed innovation law should be roughly straight. Under Gaussian
innovations one also overlays in-sample bands $\pm1.96\,\hat\sigma_t$ on $r_t$.
\end{definition}

\begin{definition}{ARMA--GARCH model}{p10-armagarch}
Because a GARCH process is itself white noise, it can serve as the innovation of
an ARMA model: with $r_t$ a $\mathrm{GARCH}(\tilde p,\tilde q)$ process,
\[
X_t=\sum_{j=1}^{p}\phi_j X_{t-j}-\sum_{k=0}^{q}\theta_k\,r_{t-k},
\qquad \theta_0=-1,
\]
where the $\phi_j,\theta_k$ satisfy the usual $\mathrm{ARMA}(p,q)$ stationarity
and invertibility conditions. This captures a conditional mean (ARMA) and a
time-varying conditional variance (GARCH) at once; one must only be careful with
the likelihood, as the errors are no longer i.i.d.\ Gaussian.
\end{definition}

% ----------------------------------------------------------------------
\subsection{Theorems \& Propositions}

\begin{proposition}{Conditional and unconditional moments of ARCH(1)}{p10-archmoments}
For the $\mathrm{ARCH}(1)$ model with $\alpha_1<1$:
\[
\E[r_t\mid\mathcal F_{t-1}]=0,\quad
\Var(r_t\mid\mathcal F_{t-1})=\sigma_t^2,\quad
\E[r_t]=0,\quad
\Var(r_t)=\frac{\alpha_0}{1-\alpha_1}.
\]
\textit{Proof idea:} $\sigma_t$ is $\mathcal F_{t-1}$-measurable and $\eps_t$ is
independent of $\mathcal F_{t-1}$, so
$\E[r_t\mid\mathcal F_{t-1}]=\sigma_t\E[\eps_t]=0$ and
$\E[r_t^2\mid\mathcal F_{t-1}]=\sigma_t^2\E[\eps_t^2]=\sigma_t^2$. Iterated
expectation gives $\E[r_t]=\E[\E[r_t\mid\mathcal F_{t-1}]]=0$ and
$\Var(r_t)=\E[\sigma_t^2]=\alpha_0+\alpha_1\Var(r_{t-1})$; under stationarity
$\Var(r_t)=\Var(r_{t-1})$, solve to get $\alpha_0/(1-\alpha_1)$.\\
\textit{Why:} establishes that ARCH is mean-zero with a positive finite
unconditional variance exactly when $\alpha_1<1$ (necessary \emph{and}, for
$\mathrm{ARCH}(1)$, sufficient for weak stationarity).
\end{proposition}

\begin{proposition}{ARCH/GARCH is white noise (uncorrelated but dependent)}{p10-whitenoise}
Any $\mathrm{ARCH}(p)$ or $\mathrm{GARCH}(p,q)$ process is mean-zero and serially
uncorrelated: for $\tau\ne0$,
\[
\E[r_t]=0,\qquad \Cov(r_{t+\tau},r_t)=0,\qquad\text{hence}\qquad \Corr(r_{t+\tau},r_t)=0 .
\]
\textit{Proof idea:} take $\tau>0$. Since $r_t$ is
$\mathcal F_{t+\tau-1}$-measurable and $\sigma_{t+\tau}$ is too, condition on
$\mathcal F_{t+\tau-1}$:
\[
\Cov(r_{t+\tau},r_t)=\E[r_{t+\tau}r_t]
=\E\!\big[\,r_t\,\sigma_{t+\tau}\,\E[\eps_{t+\tau}\mid\mathcal F_{t+\tau-1}]\big]
=\E[r_t\sigma_{t+\tau}\cdot 0]=0 .
\]
\textit{Why:} the central qualitative fact --- ARCH/GARCH leaves the
\emph{linear} correlation structure flat (white noise) while creating strong
\emph{nonlinear} dependence in the squares. The conditional-mean and
conditional-variance modelling problems decouple.
\end{proposition}

\begin{theorem}{Isserlis' theorem (Wick formula)}{p10-isserlis}
If $(X_1,\dots,X_N)$ is a zero-mean multivariate normal vector and $N$ is even,
\[
\E[X_1\cdots X_N]=\sum_{p\in P_N}\ \prod_{\{i,j\}\in p}\Cov(X_i,X_j),
\]
where $P_N$ is the set of all pairings (perfect matchings) of $\{1,\dots,N\}$.
For $N=4$,
\[
\E[X_1X_2X_3X_4]=\E[X_1X_2]\E[X_3X_4]+\E[X_1X_3]\E[X_2X_4]+\E[X_1X_4]\E[X_2X_3].
\]
\textit{Proof idea:} expand the Gaussian moment generating function / cumulant
expansion; only pair-cumulants (covariances) survive for a Gaussian.\\
\textit{Why:} gives $\E[\eps_t^4]=3(\E[\eps_t^2])^2=3$ for standard Gaussian
innovations, the key input to the ARCH fourth-moment and kurtosis computations
below.
\end{theorem}

\begin{proposition}{Fourth moment and kurtosis of ARCH(1)}{p10-kurtosis}
For $\mathrm{ARCH}(1)$ with Gaussian innovations, the unconditional fourth moment
is finite iff $0\le\alpha_1^2<\tfrac13$, in which case
\[
\E[r_t^4]=\frac{3\alpha_0^2(1+\alpha_1)}{(1-\alpha_1)(1-3\alpha_1^2)},
\qquad
\kappa=\frac{\E[r_t^4]}{\Var(r_t)^2}=\frac{3\,(1-\alpha_1^2)}{1-3\alpha_1^2}\ >\ 3
\quad(\alpha_1>0).
\]
\textit{Proof idea:} $\E[r_t^4\mid\mathcal F_{t-1}]
=\sigma_t^4\,\E[\eps_t^4\mid\mathcal F_{t-1}]=3\sigma_t^4=3(\alpha_0+\alpha_1 r_{t-1}^2)^2$
(Isserlis). Take expectations, substitute $\E[r_{t-1}^2]=\alpha_0/(1-\alpha_1)$
and $\E[r_{t-1}^4]=\E[r_t^4]$ by stationarity, and solve the linear equation;
divide by $\Var(r_t)^2=\alpha_0^2/(1-\alpha_1)^2$ for $\kappa$.\\
\textit{Why:} shows ARCH \emph{generates} excess kurtosis (fat tails) endogenously,
even with Gaussian $\eps_t$ --- matching the heavy tails seen in returns.
\end{proposition}

\begin{proposition}{ACF of the squares of an ARCH(1)}{p10-acfsq}
For $\mathrm{ARCH}(1)$ with $0<\alpha_0$ and $0<\alpha_1<1/\sqrt3$,
\[
\Corr(r_{t+\tau}^2,r_t^2)=\alpha_1^{\abs\tau},\qquad \tau\in\Ints,
\]
i.e.\ the autocorrelation of the squared series decays \emph{geometrically}.\\
\textit{Proof idea:} condition on $\mathcal F_{t+\tau-1}$ to peel one step:
$\E[r_{t+\tau}^2 r_t^2\mid\mathcal F_{t+\tau-1}]=r_t^2\sigma_{t+\tau}^2
=r_t^2(\alpha_0+\alpha_1 r_{t+\tau-1}^2)$. Taking expectations and subtracting
$(\E[r_t^2])^2$ yields the recursion
$\Cov(r_{t+\tau}^2,r_t^2)=\alpha_1\Cov(r_{t+\tau-1}^2,r_t^2)$ (the constant terms
cancel using $\E[r_t^2]=\alpha_0/(1-\alpha_1)$). Iterating gives
$\Cov(r_{t+\tau}^2,r_t^2)=\alpha_1^{\tau}\Var(r_t^2)$, hence the stated
correlation.\\
\textit{Why:} explains the slowly decaying ACF of squared returns and warns that
ARCH(1) imposes \emph{purely geometric} memory in volatility --- often too fast
for real data, motivating GARCH.
\end{proposition}

\begin{proposition}{Stationarity \& variance of GARCH(p,q)}{p10-garchstat}
A $\mathrm{GARCH}(p,q)$ process has a weakly stationary solution \textbf{iff}
\[
\sum_{j=1}^{p}\alpha_j+\sum_{j=1}^{q}\beta_j<1,
\]
and then its unconditional variance is
\[
\Var(r_t)=\frac{\alpha_0}{1-\sum_{j=1}^{p}\alpha_j-\sum_{j=1}^{q}\beta_j}.
\]
\textit{Proof idea:} take unconditional expectations of
$\sigma_t^2=\alpha_0+\sum_j\alpha_j r_{t-j}^2+\sum_j\beta_j\sigma_{t-j}^2$ using
$\E[r_{t-j}^2]=\E[\sigma_{t-j}^2]=\Var(r_t)$ at stationarity, then solve; positivity
of the denominator forces the stated bound.\\
\textit{Why:} the master condition for fitting GARCH --- it must hold for the
fitted parameters, and the closed-form variance gives the long-run volatility
level. (More: $\sum\alpha+\sum\beta<1$ is sufficient for strict stationarity; for
some noise laws including Gaussian, the boundary $\sum\alpha+\sum\beta=1$ still
admits a strictly --- but not weakly --- stationary solution: the IGARCH case.)
\end{proposition}

\begin{proposition}{Fourth-moment condition for GARCH(1,1)}{p10-garch4}
For a $\mathrm{GARCH}(1,1)$ with Gaussian noise, the unconditional fourth moment
$\E[r_t^4]$ exists \textbf{iff}
\[
3\alpha_1^2+2\alpha_1\beta_1+\beta_1^2<1 .
\]
\textit{Proof idea:} square the recursion for $\sigma_t^2$, take expectations
using $\E[\eps_t^4]=3$ (Isserlis) and stationarity of $\E[\sigma_t^4]$; the
coefficient of $\E[\sigma_{t}^4]$ on the right is $3\alpha_1^2+2\alpha_1\beta_1+\beta_1^2$,
which must be $<1$ for a finite solution.\\
\textit{Why:} the analogue of $\alpha_1^2<1/3$; needed before quoting any
formula for the kurtosis or the ACF of squared GARCH returns.
\end{proposition}

\begin{proposition}{Maximum likelihood estimation (conditional Gaussian)}{p10-mle}
Conditioning on the first observation $r_1$ to avoid specifying a marginal law
for the initial value, the conditional likelihood factorises as
\[
L_c(\theta\mid r_1,\dots,r_n)=\prod_{t=2}^{n}f(r_t\mid r_{t-1},\dots,r_1;\theta).
\]
Under Gaussian innovations the $\mathrm{ARCH}(1)$ model gives
$r_t\mid\mathcal F_{t-1}\sim\mathcal N(0,\ \alpha_0+\alpha_1 r_{t-1}^2)$, so
\[
L_c(\alpha_0,\alpha_1)=\prod_{t=2}^{n}\frac{1}{\sqrt{2\pi(\alpha_0+\alpha_1 r_{t-1}^2)}}
\exp\!\left(-\frac{r_t^2}{2(\alpha_0+\alpha_1 r_{t-1}^2)}\right),
\]
maximised over $\alpha_0>0,\ \alpha_1\ge0$ (numerically; GARCH is analogous with
$\sigma_t^2$ computed by its recursion).\\
\textit{Proof idea:} the chain rule for densities plus conditional normality of
each $r_t$ given the past.\\
\textit{Why:} the standard estimation route for ARCH/GARCH; the per-term variance
$\sigma_t^2$ is what is being fitted, not a constant $\sigma^2$.
\end{proposition}

% ----------------------------------------------------------------------
% ----------------------------------------------------------------------
